\section{Introduction}\label{sec:intro}

Large language models routinely produce \emph{competent} text but rarely produce text whose \emph{interpretive multiplicity} is deliberately authored. Where humans use a single utterance to gesture toward several incompatible readings simultaneously --- the Wittgensteinian \emph{aspect-shift}~\citep{WittgensteinPI} --- LLMs typically collapse onto a single dominant reading dictated by the prior. We argue that the missing structural element is not scale or alignment but a \emph{recursive self-reflexivity layer}: an explicit operator that re-examines its own output against the same constraint manifold that produced it, and is permitted to revise the surface when the alternatives are not yet visible in it.

Abhinavagupta's Pratyabhij\~n\=a (``recognition'') school of Kashmir \'Saiva phenomenology developed an unusually fine-grained taxonomy of exactly this layer~\citep{Lawrence2021RewritingPratyabhijna,RatieFlamandLeavitt2025}. The five-\iast{\'sakti} cascade (\iast{cit}, \iast{\=ananda}, \iast{icch\=a}, \iast{j\~n\=ana}, \iast{kriy\=a}) describes the trajectory by which an undifferentiated luminous ground (\iast{prak\=a\'sa}) becomes a determinate utterance, and the \iast{vimar\'sa} layer is the moment of self-touching --- the recognition that the utterance was authored by the very same ground that contemplates it. \iast{Apohana-\'sakti} (the contrastive exclusion of all-other-than-this) is precisely the geometric move that Bayesian model reduction~\citep{Friston2018BayesianModelReductionArxiv,FristonPenny2011} formalises as posterior switch under a parameter-pruning prior.

Our contribution is to make this isomorphism concrete and testable. We construct the \emph{Pratyabhij\~n\=a Creative Engine} (PCE), a Claude Code plugin where each \iast{\'sakti} is a real computation operating on a real local LLM substrate, and where the \iast{vimar\'sa} layer is a measurable recursive trigger gated on aspect-multiplicity, novelty, and free-energy reduction. We then evaluate PCE on four domains where the recursive layer should help if the theory is right: Wittgenstein-style poetry interpretation (where two readings must coexist), poetry generation under the POEMetric advanced-creative-abilities axes~\citep{Ye2025POEMetric}, the alternative-uses task component of CreativityPrism~\citep{CreativityPrism2025}, and a scientific-creativity slice of BIG-Bench Hard~\citep{Suzgun2022BBH,Tian2024MacGyverBench}.

The rest of the paper is organised as follows. \S\ref{sec:related} reviews prior work coupling Indic philosophical architectures to active-inference / LLM systems. \S\ref{sec:pratyabhijna_background} specifies the Pratyabhij\~n\=a primitives we operationalise. \S\ref{sec:active_inference_background} states the active-inference / BMR primitives. \S\ref{sec:engine} defines the seven operators and the consolidation cycle. \S\ref{sec:plugin} describes the plugin surface. \S\ref{sec:methods} formalises the experimental design and the pre-registered hypotheses. \S\ref{sec:benchmarks} describes the items and scoring. \S\ref{sec:results} reports the paired contrasts; \S\ref{sec:discussion}, \S\ref{sec:limitations}, and \S\ref{sec:conclusion} discuss interpretation, threats to validity, and the negative-result obligation.
